% ============================================================================
% valeurs.tex -- SOURCE UNIQUE des resultats mesures cites dans le texte
% ============================================================================
%
% Pourquoi ce fichier. Le memoire et l'article citent les memes chiffres. Une
% remesure corrigeait l'un et oubliait l'autre : la valeur 80,8 a ainsi ete
% corrigee dans le memoire en gardant 77,2 dans l'article, et il a fallu une
% relecture externe pour s'en apercevoir. Les macros ci-dessous suppriment la
% possibilite meme de cette divergence : les deux documents lisent le meme
% fichier, et une remesure ne se met a jour qu'ICI.
%
% PORTEE. Ces macros couvrent les valeurs CITEES DANS LE TEXTE, pas les
% cellules des tableaux : un tableau est un releve, il se regenere en entier
% depuis son journal quand le banc est rejoue. C'est le texte qui se perime
% sans qu'on le voie, parce que rien ne le relie a la mesure.
%
% REGLE. N'ajouter ici que ce qui est cite dans les DEUX documents, ou qui a
% deja bouge une fois. Le reste alourdirait sans proteger.
% ============================================================================

% -- ecarts garantis de reference, configuration courante, plafond 300 ------
\newcommand{\ecartVingt}{80{,}8}      % n=20, corr=0, graine 1
\newcommand{\ecartTrente}{78{,}6}     % n=30, corr=0, graine 1
\newcommand{\ecartQuarante}{96{,}5}   % n=40, corr=0, graine 1

% -- les memes lignes dans la configuration d'origine -----------------------
\newcommand{\ecartVingtOrig}{98{,}6}
\newcommand{\ecartTrenteOrig}{93{,}9}
\newcommand{\ecartQuaranteOrig}{96{,}6}

% -- archive comme livrable : medianes sur huit instances -------------------
\newcommand{\cardMed}{34{,}7}         % cardinalite, % de Z(E)
\newcommand{\hvMed}{91{,}5}           % hypervolume, % de Z(E)
\newcommand{\epsMed}{0{,}137}         % epsilon+ normalise
\newcommand{\puretePts}{102}          % points sur lesquels la purete vaut 1

% -- seconde route et forme unifiee ----------------------------------------
\newcommand{\gainGeomMax}{38{,}77}    % meilleur gain, seconde route
\newcommand{\gainUnifMax}{88{,}11}    % resserrement max, route a un tour

% -- sensibilite : ligne de production sur la strate B ---------------------
\newcommand{\ecartProd}{12{,}14}

% -- encadrements verifies -------------------------------------------------
\newcommand{\nbEncadrements}{608}     % 320 au lot cible + 288 au lot fige

% -- cout en temps de beta = 80 et rho = 4 (bench_temps.py) -----------------
\newcommand{\bruitTempsA}{20{,}6}     % temoin, strate A : ecart max a 1
\newcommand{\bruitTempsB}{8{,}0}      % temoin, strate B
\newcommand{\coutTemps}{1{,}043}      % rapport median b=80 r=4, strate B
\newcommand{\coutTemoin}{1{,}018}     % le temoin sur la meme strate

% -- les deux facteurs de la remarque « mediane bimodale », strate B --------
\newcommand{\gainMediane}{11{,}74}    % lns constante (0,4 ; 0,4)

% -- profondeur des chaines de reparation (bench_voisinage.py) -------------
\newcommand{\nbReparations}{731}      % total des chaines observees
\newcommand{\profZero}{30{,}2}        % % de chaines de profondeur 0
\newcommand{\profUn}{46{,}6}          % % de chaines a un pas
